\chapter{Background and Related Work}\label{ch:background}
This chapter introduces the foundational concepts of systems of systems (\secref{sec:background-sos}), digital twins (\secref{sec:background-dt}), and event-driven and complex event processing (\secref{sec:background-cep}), and key ideas of dependability in systems (\secref{sec:background-dependability}).

% -------------------------------------------------------------
\section{Systems of Systems}\label{sec:background-sos}
A System of Systems (SoS) is defined by \citet{maier1998architecting} as a collection of operationally and managerially independent systems that collaborate to achieve outcomes beyond their individual capabilities. Operational independence implies that each constituent system can function autonomously and continue to provide useful services outside the SoS. Managerial independence reflects that constituent systems are developed, owned, and governed by different authorities, coordinating voluntarily rather than through centralized control. Unlike monolithic systems, an SoS is characterized by decentralization, autonomy, and continuous evolution. Participation is dynamic: constituent systems may join, leave, or fail at any time, while the overall SoS continues to operate through partial composition. Coordination is achieved through negotiated interactions among heterogeneous systems pursuing distinct objectives. As a result, collective behaviour emerges from system interactions, enabling capabilities that no individual constituent can produce in isolation.

Subsequent work has refined and expanded this perspective. \citet{boardman2006system} emphasized SoS as purposeful organizations of autonomous systems united by overlapping goals and decentralized governance. 
Building on these foundations, \citet{nielsen2015systems} synthesized prior characterizations into a multi-dimensional framework for analyzing SoS. Within this framework, constituent systems exhibit a set of characteristic properties that shape their collective behaviour. \textit{Autonomy} refers to the extent to which a constituent system governs its behaviour according to its own internal goals rather than directives imposed by the SoS. \textit{Independence} captures the ability of a constituent system to operate and deliver functionality even when detached from the SoS. \textit{Distribution} reflects the spatial and logical separation of constituent systems, requiring coordination across dispersed and concurrent elements. \textit{Evolution} accounts for long-term change in system composition, functionality, and structure, while \textit{reconfiguration} refers to the ability of the SoS to adapt its composition and connectivity in real time during operation. \textit{Emergence} represents behaviours that arise from interactions among constituent systems and may not be predictable from the design of individual constituents. \textit{Interdependence} captures the mutual reliance between systems in achieving shared objectives, despite their operational independence. Finally, \textit{interoperability} describes the ability of heterogeneous systems to exchange data and services in order to function cohesively within the SoS.

Over time, the SoS paradigm expanded beyond defense into domains such as transportation, energy, emergency management, and digital commerce. With each demonstrating how loosely coupled systems can achieve stability through adaptive coordination. In transportation, air traffic and maritime networks exemplify distributed SoS, where autonomous subsystems collaborate through shared information to balance efficiency and safety. Smart energy grids integrate independent generators, utilities, and consumers into adaptive infrastructures that stabilize demand and supply in real time. Emergency management systems combine autonomous agencies and sensor networks to enable coordinated response under unpredictable conditions. Even digital platforms such as e-commerce ecosystems function as socio-technical SoS, linking payment systems, logistics providers, and supply chains into data-driven marketplaces \cite{nielsen2015systems}.


Over time, the SoS paradigm extended beyond defense into domains such as transportation, energy, emergency management, and digital commerce \cite{nielsen2015systems}. In transportation, air traffic management and maritime networks are commonly cited examples of distributed SoS, where autonomous subsystems exchange information to support traffic flow and safety management. Smart energy grids integrate independent generators, utilities, and consumers into coordinated infrastructures that balance electricity supply and demand in real time. Emergency management systems combine autonomous agencies and sensor networks to support coordinated response during disruptive events. Digital platforms such as e-commerce ecosystems similarly exhibit SoS characteristics by linking payment systems, logistics providers, and supply chains into interconnected socio-technical environments \cite{nielsen2015systems}.



% -------------------------------------------------------------
\section{Digital Twins}\label{sec:background-dt}
A Digital Twin (DT) is a high-fidelity virtual representation of a physical asset or system that remains continuously synchronized with its real-world counterpart through bi-directional data exchange and feedback control \cite{david2024infonomics}. DTs extend cyber-physical integration by enabling physical and digital entities to co-evolve: data from the physical system updates the virtual model, while the model can influence the physical system through predictive or corrective actions \cite{tao2018digital}. This synchronization, supported by IoT-enabled sensing and networked infrastructures, allows the DT to replicate both the structure and dynamic behaviour of its physical twin in real time.

A DT comprises three primary components \cite{vanderhorn2021digital}:
(i) the \emph{physical system}, representing the asset, process, or
environment being observed; (ii) the \emph{virtual representation},
which abstracts this reality using computational models; and (iii) the
\emph{interconnections} that enable continuous information exchange
between the physical and digital components. These interconnections support physical-to-virtual flows that update the digital state based on observations, as well as virtual-to-physical flows that communicate predictions, recommendations, or control signals back to the physical system.

To clarify related forms of digital representation, \citet{kritzinger2018digital} distinguish between a digital model, a digital shadow, and a digital twin. A digital model represents a static computational representation of a physical system without live data connectivity, while a digital shadow supports one-way updates from the physical system to the digital representation. A digital twin, in contrast, maintains a recurring, bidirectional exchange that enables continuous synchronization and supports operational services.

% Event driven lifecycle
DTs operate according to an event-driven lifecycle in which synchronization is triggered by changes observed in the physical system. These changes are communicated as discrete, time-stamped observations that initiate state updates within the digital representation. More generally, DT synchronization behaviour is characterized in terms of discrete triggers that determine when such updates occur, including state changes, timed triggers, contextual conditions, and data-dependent thresholds \cite{alghamdi2024synchronization}. The updated digital state may then be used for simulation or analysis, followed by optimization or control actions that influence physical operation. As this lifecycle repeats, DT behaviour becomes dependent on the timing, ordering, and continuity of observed events. Consequently, DT synchronization is governed by event semantics that determine when and under what conditions digital state updates occur. These semantics are also formalized in established DT standards. ISO~23247 defines event-based synchronization as a mechanism for maintaining consistency between a digital twin and its physical system through event-triggered updates \cite{iso23247_1_2021}. Asset Administration Shell–based architectures similarly specify explicit event constructs and their use for asynchronous communication between DTs and external systems \cite{aas_details_2018_p1}.


% -------------------------------------------------------------
% \section{Event-Driven Processing and Complex Event Processing}\label{sec:background-cep}
% Event-driven processing is a computational model in which system behaviour advances in response to the occurrence of events \cite{kommera2020power}. An event represents a discrete, observable occurrence within a system or its environment, such as a state change, measurement, or notable condition, and is commonly modeled as a record carrying a payload and a timestamp \cite{etzion2010event}. Events are produced asynchronously by sources such as sensors, software components, workflow engines, or external services, and are collected into event streams that reflect the evolving behaviour of a system over time. CEP extends the event-driven model by enabling the detection of higher-level situations through the interpretation and correlation of multiple events over time \cite{cugola2015complex}. CEP systems evaluate continuous event streams against declarative rules or patterns that specify temporal, logical, or causal relationships among events, including sequencing, aggregation, concurrency, negation, and time windows \cite{cugola2012processing,cugola2015complex}. When a pattern is satisfied, the system emits a composite or complex event representing the detection of a meaningful situation.

% To support incremental pattern detection, CEP engines maintain internal state to track partial matches over time. Common execution models rely on automata-based or state-machine-based techniques to manage long-running patterns and temporal constraints, introducing trade-offs between expressiveness, memory usage, and performance \cite{cugola2015complex}. CEP originated in academic systems such as SASE and Cayuga, which introduced declarative models for event sequence and pattern detection over unbounded streams \cite{wu2006high,brenna2007cayuga}. These ideas were later adopted in industrial CEP engines such as Esper and Siddhi for real-time event processing \cite{esper_reference,suhothayan2011siddhi}, and increasingly integrated into large-scale distributed stream processing frameworks such as Apache Flink \cite{carbone2015apache,dayarathna2018recent}.

% CEP has been applied across domains including logistics, finance, manufacturing, infrastructure monitoring, and healthcare, where timely detection of evolving situations is required \cite{buchmann2009complex,manchana2021event}. As systems have evolved toward distributed, event-driven and microservice-based architectures, CEP has emerged as a foundational abstraction for interpreting and coordinating behaviour across heterogeneous and distributed event sources \cite{kommera2020power,dayarathna2018recent}.


% -------------------------------------------------------------
\section{Dependability and Reliability in Systems}
\label{sec:background-dependability}

Dependability provides a foundational framework for reasoning about the trustworthiness of computing systems. \citet{laprie1985dependable} define dependability as the quality of delivered service such that reliance can justifiably be placed on it, a view later formalized by \citet{avizienis2001fundamental,avizienis2004basic} as a system-level property describing the ability to deliver correct service in the presence of faults and adverse conditions. Within this framework, a fault represents the underlying cause of an impairment, an error corresponds to an incorrect internal system state, and a failure occurs when an error results in an externally observable deviation from correct service \cite{avizienis2004basic}. 

Dependability is characterized through a set of attributes, including availability, reliability, safety, integrity, and maintainability \cite{avizienis2004basic}. Availability refers to the readiness of a system to deliver correct service when required, whereas reliability describes the continuity of correct service over time. These attributes are complementary: reliability quantifies the probability of failure-free operation over a specified time interval, while availability captures the proportion of time a system is operational, accounting for both failures and restorations \cite{laprie1985dependable}. Safety concerns the avoidance of catastrophic consequences arising from system failures, integrity ensures that system state and data are protected from improper modification, and maintainability reflects the ease with which a system can be repaired or adapted.

Dependability theory further distinguishes four complementary means for achieving dependable systems: fault prevention, fault removal, fault forecasting, and fault tolerance \cite{avizienis2004basic}. Fault prevention and fault removal aim to reduce the introduction and persistence of faults through design, verification, testing, and maintenance, while fault forecasting estimates the likelihood and consequences of faults and failures. Fault tolerance addresses the presence of residual faults by enabling systems to continue delivering acceptable service during operation. Fault-tolerant systems rely on mechanisms for fault detection, containment, and recovery, often supported by redundancy in hardware, software, information, or computation, and realized through architectural structures that isolate faults, prevent error propagation, and restore system state through recovery or reconfiguration \cite{nelson2002fault,hanmer2013patterns}.


% -------------------------------------------------------------
% Failures may propagate across component dependencies, as components inherit the failure semantics of those they rely on, unless such failures are explicitly masked \cite{cristian1991understanding}. While system-level dependability cannot be reduced to the reliability of individual components alone, prior work emphasizes that improving the dependability and failure semantics of critical components is a necessary condition for achieving trustworthy overall system behaviour \cite{littlewood2000software,hanmer2013patterns}.





% Having trouble integrating in contract based composition
% Classical contract-based design (Benveniste et al.), addresses the composition of systems under strong assumptions about component behaviour, interfaces, and integration. - relies on fixed environments and well-defined assumptions and guarantees to ensure correctness under composition. 
% SoTS however, do not operate under fixed environments: constituent systems may join or leave, interaction patterns evolve over time, 
% So strong contract-based composition cannot be relied on in these environments
% This thesis addresses the complementary problem of supporting reliable event-driven computation in systems where composition is inherently weak and dynamic.



